\section{Problem Setting and Definitions}
The goal of reward design is to return a shaped reward function for a ground-truth reward function that may be difficult to optimize directly (e.g., sparse rewards); this ground-truth reward function may only be accessed via queries by the designer.  We first introduce the formal definition from~\citet{singh2010where}, which we then adapt to the program synthesis setting, which we call \textit{reward generation}. 

\begin{definition} 
\label{def:rdp}
(Reward Design Problem~\citep{singh2010where}) A \emph{reward design problem} (RDP) is a tuple $P = \langle M, \mathcal{R}, \pi_M, F \rangle$, where $M = (S,A,T)$ is the \textit{world model} with state space $S$, action space $A$, and transition function $T$. $\mathcal{R}$ is the space of reward functions; $\mathcal{A}_M(\cdot): \mathcal{R} \rightarrow \Pi$ is a learning algorithm that outputs a policy $\pi:S \rightarrow \Delta(A)$ that optimizes reward $R \in \mathcal{R}$ in the resulting \textit{Markov Decision Process} (MDP), $(M,R)$; $F: \Pi \rightarrow \mathbb{R}$ is the \textit{fitness} function that produces a scalar evaluation of any policy, which may only be accessed via policy queries (i.e., evaluate the policy using the ground truth reward function). In an RDP, the goal is to output a reward function $R \in \mathcal{R}$ 
such that the policy $\pi:=\mathcal{A}_M(R)$ that optimizes $R$ achieves the highest fitness score $F(\pi)$.
\end{definition}


\textbf{Reward Generation Problem.} In our problem setting, every component within a RDP is specified via code. 
Then, given a string $l$ that specifies the task, the objective of the reward generation problem is to output a reward function code $R$ such that $F(\mathcal{A}_{M}(R))$ is maximized. 








